% Copyright (c) 2022 by Lars Spreng
% This work is licensed under the Creative Commons Attribution 4.0 International License. 
% To view a copy of this license, visit http://creativecommons.org/licenses/by/4.0/ or send a letter to Creative Commons, PO Box 1866, Mountain View, CA 94042, USA.

%~~~~~~~~~~~~~~~~~~~~~~~~~~~~~~~~~~~~~~~~~~~~~~~~~~~~~~~~~~~~~~~~~~~~~~~~~~~~~~
% You can add your packages and commands to the loadslides.tex file. 
% The files in the folder "styles" can be modified to change the layout and design of your slides.
% I have included examples on how to use the template below. 
% Some of these examples are taken from the Metropolis template.
%~~~~~~~~~~~~~~~~~~~~~~~~~~~~~~~~~~~~~~~~~~~~~~~~~~~~~~~~~~~~~~~~~~~~~~~~~~~~~~


\documentclass[
11pt,notheorems,hyperref={pdfauthor=whatever}
]{beamer}

\input{loadslides.tex} % Loads packages and some defined commands

\title[
% Text entered here will appear in the bottom middle
]{Project Overview}

\subtitle{\textbf{Air Piano}}

\author[
% Text entered here will appear in the bottom left corner
]{
    N. Bacon, M. Tang, A. Taylor
}

\begin{document}

% Generate title page
{
\setbeamertemplate{footline}{} 
\begin{frame}
  \titlepage
\end{frame}
}
\addtocounter{framenumber}{-1}

\section{Problem Statement and Success Criteria}
\begin{frame}

\begin{block}{Project Description}
\begin{itemize}
    \item Our project is the \textbf{Air Piano}, a
    contact-less digital instrument. Users play
    notes by making motions with their hand
    in front of a break beam sensor. The pitch of each
    note is determined by the distance from the
    device to their hand as measured by an ultrasonic
    distance sensor. Like a classical piano, pitches
    are quantized based on the standard 12-TET tuning
    system. To increase playability, we chose to limit
    the device to eight notes, forming a C major scale.
    Finally, the device includes a real-time LED
    display indicating which pitch is currently
    to be played based on the distance sensor.
\end{itemize}
\end{block}

\begin{block}{Success Criteria}
\begin{enumerate}
    \item The device plays a single note for a short
        time (configurable) on the buzzer
        only when a hand moves downwards in front of it.
    \item The pitch of the note is lower when the hand
        is closer in horizontal distance to the device
        and higher when it is further away. The pitches
        played by the device will range from C4 to C5.
    \item The LED display on the device correctly
    displays the note that would be played should the
    user lower the hand closest to the device at that
    moment.
    \item The timing requirements for the device are
        met.
\end{enumerate}
\end{block}

\end{frame}

\section{Hardware}
\begin{frame}

\begin{block}{SBC Choice}
\begin{itemize}
    \item Beaglebone Black
\end{itemize}
\end{block}

\begin{block}{Electronics Kit Components}
\begin{itemize}
    \item Breadboard
    \item Jumper wires
    \item Passive buzzer
    \item LEDs
\end{itemize}
\end{block}

\begin{block}{Add-On Sensors}
\begin{itemize}
    \item IR break beam sensor
    \item HC-SR04 ultrasonic distance sensor
\end{itemize}
\end{block}

\end{frame}

\section{Requirements}
\begin{frame}

\begin{block}{Measurable Timing Requirements}
\begin{enumerate}
    \item On note-down event as measured by the break
        beam sensor, audio output must begin
        with 10ms worst-case.
    \item On note-change event as measured by the
        distance sensor, the LED display must begin
        updating within 100ms worst-case.
    \item Break beam sensor is to be sampled at 100Hz
        with jitter $\pm$ 1ms.
    \item Distance sensor is to be sampled at 30Hz
        with jitter $\pm$ 1ms.
\end{enumerate}
\begin{itemize}
    \item \emph{See \cite{wessel_wright_2002}}
\end{itemize}
\end{block}

\end{frame}

\begin{frame}

\begin{block}{Periodic Tasks}
\begin{enumerate}
    \item Take break beam measurements every 10ms (100Hz)
    \item Take distance measurements every 33ms (30Hz)
    \item Update LED display every 50ms (20Hz)
    \item Write logs every 100 ms (10 Hz)
\end{enumerate}
\end{block}

\begin{block}{Asynchronous Task}
\begin{itemize}
    \item Upon beam break detection, play current
        note within 10ms.
\end{itemize}
\end{block}

\end{frame}

\section{Architecture}
\begin{frame}

\begin{block}{Task Architecture}
\begin{table}
\small
\begin{tabular}{@{} lll @{}}
\toprule
\textbf{Task Name} & \textbf{Priority} & \textbf{Responsibilities}\\
\midrule
Beam Sensor & 90 & Detects break beam being broken and
    plays the note determined by distance \\
    &&when triggered.\\
Distance Sensor & 70 & Measures distance from sensor and
    calculates corresponding note.\\
Display Updater & 60 & Continuously displays the current note
    indicated by the distance sensor.\\
Console & 40 & Reads special commands from console input.\\
Logger & 30 & Outputs both real-time updates to console output
    as well as logs to file.\\
\bottomrule
\end{tabular}
\end{table}
\end{block}

\begin{block}{Task Communication}
\begin{itemize}
    \item The Distance Sensor will continuously update
        a shared value with the current note that
        corresponds to the distance measurement. This
        allows the Display Updater and the Beam Sensor
        to read this shared value to quickly
        play or display the note. The data will be
        protected by a lock.
    \item Tasks will communicate with the Logger via
        buffers protected by a lock. The logger will
        only hold onto the lock while copying data to
        its own memory.
    \item The safe state will be protected by
        a lock, and different tasks will update it based
        on their specific failure-detection criteria.
\end{itemize}
\end{block}

\end{frame}
\begin{frame}

\begin{block}{Safe-State Behavior}
\begin{itemize}
    \item The safe-state behavior of the Air Piano
        will turn off all LEDs and output no sound.
        This will signal that something is wrong and
        prevent annoyance from an indefinitely
        held note.
    \item If more than 10 beam breaks are detected
        within one
        second, the safe state will be entered, as this
        indicates a possible sensor or logical
        malfunction.
    \item If the distance sensor detects a value
        outside of the playable note range, no audio
        will be produced, but the device will not enter
        the safe state. If the distance sensor value is
        clearly outside of the sensor's own range,
        the safe state will be entered.
    \item To exit the safe state, an admin will enter
        a command on the console.
    \begin{itemize}
        \item Admins may also manually enter the safe
        state via console command.
    \end{itemize}

\end{itemize}
\end{block}

\end{frame}

\section{Tests and Measurements}
\begin{frame}

\begin{block}{Timing Data}
\begin{enumerate}
    \item Capture timestamps for both beam breaks
        and audio playback to ensure deadlines are
        met. Log to file for review and testing.
    \item Also record captured distance with notes
        played to verify sensor logic and overall
        functionality. Store in file for review and
        testing.
\end{enumerate}
\end{block}

\begin{block}{Boundary and Load Cases}
\begin{enumerate}
    \item Test with high sensor load (8 beam break
        events per second) to verify deadlines are met
        and functionality is retained.
    \item Test with multiple objects around the sensor
        such as people walking in the background to
        ensure functionality isn't impeded by unwanted
        interference (no spurious note playback
        events).
    % \item Simulate slow tasks with artificial
    %     delays, making sure the rest of the system
    %     functions as described. The slowing of low-priority
    %     tasks should not interfere with high-priority
    %     tasks.
    \item Test with high CPU load by running other
        processes in the background.
\end{enumerate}
\end{block}

\end{frame}

\begin{frame}

\begin{block}{Base Timing Measurements}
\begin{itemize}
    \item While most deadlines were met, our measurements
    showed that the system occasionally missed deadlines,
    even for the highest priority tasks.
    \item Our measurements indicate very low latency
    to output the note once detected, with worst case
    latency less than 1ms. Likewise, the latency for
    the display was always less than 100ms.
\end{itemize}
\end{block}

\begin{block}{Load Conditions}
\begin{itemize}
    \item By simulating 8 beam break events per second,
    we confirmed that the system functioned as normal.
    \item From ad-hoc experiments, we determined that
    noise from the external environment can
    affect the operation of the device (e.g. objects
    near the distance sensor).
    \item At 100\% CPU load, we found our timing data
        to be comparable to the baseline.
\end{itemize}
\end{block}

\end{frame}

\begin{frame}[allowframebreaks]{References}
    \printbibliography
\end{frame}

\end{document}